% BHU PYQ -- Manually transcribed & error-corrected
% Paper: BCHF-326: Financial Markets in India
% Exam: B.Com. (Hons.) Semester VI Examination, 2015-16

\documentclass[11pt,a4paper]{article}
\usepackage[a4paper,top=2.5cm,bottom=2.5cm,left=2.8cm,right=2.8cm,headheight=14pt]{geometry}
\usepackage{amsmath,amssymb}
\usepackage[T1]{fontenc}
\usepackage[utf8]{inputenc}
\usepackage{lmodern,microtype,fancyhdr,enumitem,hyperref,lastpage,parskip}

\hypersetup{
    colorlinks=true,
    urlcolor=black,
    linkcolor=black,
    pdftitle={BCHF-326: Financial Markets in India -- BHU B.Com. (Hons.) Sem VI 2015-16}
}

\pagestyle{fancy}
\fancyhf{}
\renewcommand{\headrulewidth}{0.4pt}
\renewcommand{\footrulewidth}{0.4pt}
\fancyhead[L]{\small   \textit{Banaras Hindu University}}
\fancyhead[R]{\small   \textit{BCHF-326: Financial Markets in India}}
\fancyfoot[C]{\small Page  \thepage\ of \pageref{LastPage}}


\newlist{parts}{enumerate}{2}
\setlist[parts,1]{label=(\alph*),leftmargin=*,itemsep=5pt,topsep=3pt}
\setlist[parts,2]{label=(\roman*),leftmargin=*,itemsep=3pt,topsep=2pt}

\newcommand{\pts}[1]{\hfill   \textit{[#1]}}


\begin{document}


\begin{center}
  {\large\bfseries BANARAS HINDU UNIVERSITY}\\[2pt]
  {\normalsize B.Com. (Hons.) --- Semester VI Examination, 2015-16}\\[6pt]
  \rule{0.8\linewidth}{0.5pt}\\[6pt]
  {\large\bfseries Elective Group --- Finance}\\[4pt]
  {\large\bfseries Paper No. BCHF-326: Financial Markets in India}\\[4pt]
  {\normalsize Time: 3 Hours\hfill Maximum Marks: 70}
\end{center}

\rule{\linewidth}{0.4pt}

\smallskip



\noindent   \textit{Instructions: Answer all questions. All questions carry equal marks (14 marks each).}

\bigskip




\noindent   \textbf{Question 1.} \\
What do you understand by Industrial Securities Market? Discuss its main features.\pts{14}

\centerline{   \textbf{OR}}

What is the Long-term Loans Market? Discuss the role of the mortgages market in the Long-term Loans Market.\pts{14}

\medskip\rule{\linewidth}{0.2pt}




\noindent   \textbf{Question 2.} \\
Discuss the main features of the structure of the money market in India.\pts{14}

\centerline{   \textbf{OR}}

Why is a developed money market essential for a developed or developing economy? Discuss.\pts{14}

\medskip\rule{\linewidth}{0.2pt}

\pagebreak


\begin{center}
  \framebox[\linewidth]{\parbox{0.95\linewidth}{\small   \textbf{Scanning Anomaly Note:} In the original scanned booklet, page 32 (page 3 of this exam paper) was misprinted or scanned with page 38 from the Financial Analysis (BCH-321) paper. As a result, Questions 3 and 4 of BCHF-326 are missing from the source document, and the following accounting question is printed instead. We transcribe it here to preserve the scanned source exactly.}}
\end{center}

\medskip




\noindent   \textbf{Question 3 \& 4 (Misprinted in Booklet Scan).}\\
Give a detailed format of a cash flow statement. How does it differ from a funds flow statement?\pts{14}

\centerline{   \textbf{OR}}

The Balance Sheet of Paschim Corporation Ltd. as on 31st December, 2014 and 2015 is as follows:
\pts{14}


\begin{center}
\small
\renewcommand{\arraystretch}{1.2}

\begin{tabular}{|l|c|c|l|c|c|}
\hline
   \textbf{Liabilities} &   \textbf{2014 (Rs.)} &   \textbf{2015 (Rs.)} &   \textbf{Assets} &   \textbf{2014 (Rs.)} &   \textbf{2015 (Rs.)} \\ \hline
11\% Cum. Pref. Shares & -- & 30,000 & Land \& Building & 60,000 & 50,000 \\
Equity Shares & 1,10,000 & 1,20,000 & Plant & 30,000 & 50,000 \\
Reserve & 4,000 & 4,000 & Sundry Debtors & 40,000 & 48,000 \\
Profit \& Loss A/c & 2,000 & 2,400 & Stock & 60,000 & 70,000 \\
9\% Debentures & 12,000 & 14,000 & Bank & 2,400 & 7,000 \\
Provision for Taxation & 6,000 & 8,400 & Cash & 600 & 1,000 \\
Proposed Dividend & 10,000 & 11,600 & & & \\
Current Liabilities & 49,000 & 35,600 & & & \\ \hline
   \textbf{Total} &   \textbf{1,93,000} &   \textbf{2,26,000} &   \textbf{Total} &   \textbf{1,93,000} &   \textbf{2,26,000} \\ \hline
\end{tabular}
\end{center}

You are required to prepare a schedule of changes in working capital and a statement of Flow of Funds.

\medskip\rule{\linewidth}{0.2pt}




\noindent   \textbf{Question 5.} \\
What are the guidelines of SEBI for protecting the interest of debenture holders?\pts{14}

\centerline{   \textbf{OR}}

Discuss in brief some of the important guidelines on IPOs (Initial Public Offerings) through the on-line system (E-IPO).\pts{14}

\vfill
\rule{\linewidth}{0.4pt}

\begin{center}\small
\href{https://bhu-exam.vercel.app/}{Click here to visit our website}\\[4pt]
\href{https://me-aryan.vercel.app/}{Click here to visit our contributor}\\[4pt]
\href{https://quashan-me.vercel.app/}{Click here to visit our contributor}
\end{center}

\end{document}
